\documentclass[12pt]{article}
\usepackage{morse}

\def\lm{\Large\morse}
\pagestyle{plain}
\topmargin-1cm
\textwidth15.85cm
\textheight23.35cm
\parindent0pt
\sloppy

\begin{document}

\underline{\LARGE \Morsesym\/ -- a short font report}\\[1ex]
(The telegraphic alphabet - by {\sc Udo Heyl}, January $1^{st}$, 1998)\\
\begin{center}
\fbox{\parbox{6in}{
\begin{center}
{\it Error Reports in case of UNCHANGED versions to}\\ 
   Udo Heyl,
   Stregdaer Allee 7,
   99817 Eisenach,
   Federal Republic of Germany \\
{\it or} \\
   DANTE, Deutschsprachige Anwendervereinigung TeX e.V., 
   Postfach 10 18 40,\\ 
   69008 Heidelberg, 
   Federal Republic of Germany, 
   Email: dante@dante.de 
\end{center}
}}
\end{center}

\section{What is \Morsesym ?}

It is a \LaTeX2e - package for printing Morse code signs.

\section{How to use the \Morsesym\/ package?}

First and foremost you've got to copy the following files
\begin{itemize}
\item {\sc morse10.mf, morse.def, morse.alf} and {\sc morse.num} 
      into your Metafont-directory
      \verb/(\emtex\mfinput\morse)/
\item {\sc morse.sty} into your Style-directory 
      \verb/(\emtex\texinput\morse)/ and
\item {\sc morse10.tfm} into your Tfm-directory
      \verb/(\emtex\tfm\morse)/.
\end{itemize}
Note, however, that the paths may be different in your 
\LaTeX2e ~implementation 
(Em\TeX ~for {\sc MS-DOS}, {\sf web2c} for {\sc UNIX} etc.).
\LaTeX2e ~is absolutely required, if you want to use \Morsesym, 
which doesn't run with the {\bf ancient} \LaTeX 209. \\

The example shows you the usage of this font:
\begin{verbatim}
   \documentclass[10pt]{article}
   \usepackage{morse} %%% to include morse.sty
   \begin{document} ... 
   %..........
   {\morse ..... }
   %..........
   \end{document}
\end{verbatim}
Telegraphic code will appear now in the current size
and won't change the current shape. Of course you can input
\verb/{\Huge\morse /{\it MorseText}\verb/ }/ to manage greater
Morse code signs. After the command \verb/\morse/ numbers, stops
and letters (upper and lower case) are converted into Morse code,
exceptions are given in Table 1.

\begin{center}
\begin{table}[t]
\caption{\Large\bf The Alphabet of the Telegraphy}
\vskip20pt
\begin{tabular}{||l|c|c||l|c|c||l|c|c||} 
\hline\hline
&&&&&&&&\\
Input & Output & Value&Input & Output & Value& 
                           Input & Output & Value\\
&&&&&&&&\\\hline &&&&&&&&\\
\verb=a= & \lm a & a & \verb=n= & \lm n & n & \verb=3= & \lm 3 & 3 \\
\verb=\aAcute= &\lm\aAcute & \'{a} & \verb=o= &
                                    \lm o & o & \verb=4= & \lm 4 & 4 \\
\verb=\ae= &\lm\ae & \H{a} & \verb=\oe= & \lm\oe & \H{o} & 
                                    \verb=5= &\lm 5 & 5 \\
\verb=b= & \lm b & b & \verb=p= & \lm p & p & \verb=6= & \lm 6 & 6 \\
\verb=c= & \lm c & c & \verb=q= & \lm q & q & \verb=7= & \lm 7 & 7 \\
\verb=d= & \lm d & d & \verb=r= & \lm r & r & \verb=8= & \lm 8 & 8 \\
\verb=\ch= & \lm\ch & ch & \verb=s= & \lm s & s & 
                                    \verb=9= & \lm 9 & 9 \\
\verb=e= & \lm e & e & \verb=t= & \lm t & t & \verb=0= & \lm 0 & 0 \\
\verb=\eAcute= &\lm\eAcute & \'{e} & \verb=u= & \lm u & u &
                                    \verb=;= & \lm ; & ; \\
\verb=f= & \lm f & f & \verb=\ue= & \lm \ue & \H{u} & 
                                    \verb=,= & \lm , & , \\
\verb=g= & \lm g & g & \verb=v= & \lm v & v & \verb=:= & \lm : & : \\
\verb=h= & \lm h & h & \verb=w= & \lm w & w & \verb=?= & \lm ? & ? \\
\verb=i= & \lm i & i & \verb=x= & \lm x & x & \verb=!= & \lm ! & ! \\
\verb=j= & \lm j & j & \verb=y= & \lm y & y & 
                                           \verb=\dq= & \lm \dq & \dq \\
\verb=k= & \lm k & k & \verb=z= & \lm z & z & 
                                           \verb=\sq= & \lm \sq & \sq \\
\verb=l= & \lm l & l & \verb=1= & \lm 1 & 1 & \verb=/= & \lm / & / \\
\verb=m= & \lm m & m & \verb=2= & \lm 2 & 2 & \verb=.= & \lm . & . \\
&&&&&&&&\\
\hline\hline
\end{tabular}
\end{table}
\end{center}
\vskip-1cm
\begin{footnotesize}
\begin{minipage}{7.4cm}
In the year 1832 the American painter and director of the  
Plastic Arts Academy {\sc Samuel Morse} had the idea to use electromagnetism 
in order to set up a telegraphic connection. 
After some different unsuccessful experiments he could propose a model of 
his \dq Recording electric telegraph\dq ~to the New York University in 1835.  
In Washington he took out a patent for his invention in 1837.
At the same time 
   {\sc Charles Wheatstone} in England and
   {\sc Karl August Steinheil} in Bavaria
had made up telegraphs.
The one by the last mentioned was the same on principle as {\sc Morse}'s
telegraph, 
but wasn't as useful in drawing longer lines \hfill
because its fine and complex mecha-
\end{minipage}
\hskip1cm
\begin{minipage}{7.4cm}
nism. That is why {\sc Steinheil} himself
recommended the installation of {\sc Morse}'s telegraph,
which was the better one in later experiments too because of its
plainness and easy usage up to now.
{\sc Morse}'s system consists of an armature with a pin, which is
attracted by a powerized electric magnet.
The pin prints dots or dashes onto a passing tape of paper,
in proportion as the power circuit is complete for shorter or longer
periods.
The foregoing alphabet is composed out of this dots and dashes.
The character \'{a} is an 
innovation for the Hungarian language; 
\H{a}, \H{o} and \H{u} for the German vocabulary; and in Polish words
the characters    \c{a}, \c{e} and \'{o} are described by
                  \H{a}, \'{e} and \H{o}.
\end{minipage}
\end{footnotesize}
\end{document}
